\subsection{Documentation}

The document is split into the following sections:

\begin{enumerate}
    \item Ideas
    \item Requirements
    \item Design
    \item Implementation
    \item Deployment
    \item Issue Resolution
\end{enumerate}

Note that while the process of taking an Idea all the way through to Deployment
is cyclical each Module that is taken from start to finish of the cycle will be
considered complete.

\subsubsection{Ideas}

In the ideas section we essentially have our introduction to the business problem
and the proposed business solutions.  This section should include a high level 
overview and should provide just enough information to be succinct and sufficient
to extract meaningful and actionable business requirements to be described in the
Requirements section.

This section documents the ``what if?'' of an idea.  Later sections will document
the ``why'', ``what'', ``how'', ``when'', ``did it?'', and ``is it?''.

The ``did it?'' is particularly important in that it answer the question from 
``if we do X then we will get Y?''.  This is evidence based design and should
be considered a crucial step in ensuring that features and functionality actually
achieve and produce some positive change.  If they don't, then a new solution
or idea should be discussed and implemented.

The \emph{``is it?''} is meant to categorize the unit, feature, and integration 
tests.  Note that while unit tests will be developed for custom code, we will also
have feature and integration tests for general software products.  For this 
purpose we will use the \emph{Common Test} framework of the erlang programming
language\footnote{See section \ref{sec:commontest}}.

\subsubsection{Requirements}
Once the general business problem and proposed solution have been documented, we 
must extract formal requirements that will be used to deliver a functional solution.

Note that we are not limiting ourselves to programming solutions, but rather, 
technological solutions at the architectural level (i.e. the solution code be an off
the shelf software product, deploying a virtual server, or writing custom source code).

The requirements will be written primarily in the Agile style (As {\emph{User}} 
I want {\emph{Feature}} to achive {\emph{Functionality}}).  Here \emph{User} can 
be any role, person, or group.  It does not have to be something software specific.
Similarly, the \emph{Functionality} does not have to be a custom built piece of 
code, although it will be some IT feature, with possible some manual process to
augment it.  The \emph{Functionality} will be the goal or result that we want to 
achieve via the requested \emph{Feature}.  Note that the \emph{Feature} will 
usually be some increase in efficiency, decrease in work/cost, or improvement
to some company, group, or user.

This section formalizes the ``what'' and the ``why''.

\subsubsection{Design}

